\documentclass[11pt]{article}

\usepackage[margin=1in]{geometry}
\usepackage{amsmath}
\usepackage{booktabs}
\usepackage[tableposition=top]{caption}
\usepackage{graphicx}
\usepackage{natbib}
\bibliographystyle{plainnat}

% Keep this package last. It adds the margin link from each figure to the
% script that made it, and the commit stamp.
\usepackage{showyourwork}

\title{A Reproducible Paper}
\author{Your Name}

\begin{document}

\maketitle

\begin{abstract}
  Every figure, number, and table in this paper is rebuilt from code by
  \texttt{./build.sh}. The icon next to each figure caption links to the
  script that made the figure, at the commit that built this PDF. Replace
  this text with your paper.
\end{abstract}

\section{How this paper is built}

This paper uses showyourwork \citep{Luger2021}. A figure is never a file you
drop into the repository: it is the output of a script in
\texttt{src/scripts/}, and the build runs that script whenever the script or
its data change.

Figure~\ref{fig:walkers} is the simplest case. Its figure environment
contains \verb|\script{walkers.py}|, and showyourwork infers the rest.

\begin{figure}[h]
  \centering
  \includegraphics[width=0.6\linewidth]{figures/walkers.pdf}
  \caption{Five random walks of 200 steps.}
  \label{fig:walkers}
  \script{walkers.py}
\end{figure}

Figure~\ref{fig:spread} comes from a cached result. This is an expensive
run whose output file is committed to the repository and recomputed only by
\texttt{./build.sh --regenerate}. The script that plots it reads that file,
so the \texttt{Snakefile} lists the file as an input of the figure.

\begin{figure}[h]
  \centering
  \includegraphics[width=0.6\linewidth]{figures/spread.pdf}
  \caption{Mean squared displacement of 10\,000 random walkers. The dashed
    line is the theory, $\langle x_t^2 \rangle = t$.}
  \label{fig:spread}
  \script{spread.py}
\end{figure}

Numbers in the text come from code too. A linear fit to
Figure~\ref{fig:spread} has slope \variable{output/slope.txt}, and
Table~\ref{tab:results} lists the values at three steps. Both files are
written by \texttt{src/scripts/results.py}.

\begin{table}[h]
  \centering
  \caption{Mean squared displacement at three steps.}
  \label{tab:results}
  % showyourwork compiles ms.tex once before any script runs. A missing
  % \input is fatal there, so guard every \input of a generated file.
  \IfFileExists{output/results_table.tex}{\input{output/results_table.tex}}{}
\end{table}

\bibliography{bib}

\end{document}
